\documentclass[addpoints,spanish, 12pt,a4paper]{exam}
%\documentclass[answers, spanish, 12pt,a4paper]{exam}
\printanswers
\pointpoints{punto}{puntos}
\hpword{Puntos:}
\vpword{Puntos:}
\htword{Total}
\vtword{Total}
\hsword{Resultado:}
\hqword{Ejercicio:}
\vqword{Ejercicio:}

\usepackage[utf8]{inputenc}
\usepackage[spanish]{babel}
\usepackage{eurosym}
%\usepackage[spanish,es-lcroman, es-tabla, es-noshorthands]{babel}


\usepackage[margin=1in]{geometry}
\usepackage{amsmath,amssymb}
\usepackage{multicol}
\usepackage{yhmath}

\pointsinrightmargin % Para poner las puntuaciones a la derecha. Se puede cambiar. Si se comenta, sale a la izquierda.
\extrawidth{-2.4cm} %Un poquito más de margen por si ponemos textos largos.
\marginpointname{ \emph{\points}}

\usepackage{graphicx}

\graphicspath{{../../img/}} 

\newcommand{\class}{2º Bachillerato CIT}
\newcommand{\examdate}{\today}
\newcommand{\examnum}{Parcial 2ªEv. - Global de Análisis}
\newcommand{\tipo}{A}


\newcommand{\timelimit}{105 minutos}

\renewcommand{\solutiontitle}{\noindent\textbf{Solución:}\enspace}


\pagestyle{head}
\firstpageheader{\includegraphics[width=0.2\columnwidth]{header_left}}{\textbf{Departamento de Matemáticas\linebreak \class}\linebreak \examnum}{\includegraphics[width=0.1\columnwidth]{header_right}}
\runningheader{\class}{\examnum}{Página \thepage\ of \numpages}
\runningheadrule


\usepackage{pgf,tikz,pgfplots}
\pgfplotsset{compat=1.15}
\usepackage{mathrsfs}
\usetikzlibrary{arrows}


\begin{document}

\noindent
\begin{tabular*}{\textwidth}{l @{\extracolsep{\fill}} r @{\extracolsep{6pt}} }
\textbf{Nombre:} \makebox[3.5in]{\hrulefill} & \textbf{Fecha:}\makebox[1in]{\hrulefill} \\
 & \\
\textbf{Tiempo: \timelimit} & Tipo: \tipo 
\end{tabular*}
\rule[2ex]{\textwidth}{2pt}
Esta prueba tiene \numquestions\ ejercicios. La puntuación máxima es de \numpoints. 
La nota final de la prueba será la parte proporcional de la puntuación obtenida sobre la puntuación máxima. 

\begin{center}


\addpoints
 %\gradetable[h][questions]
	\pointtable[h][questions]
\end{center}

\noindent
\rule[2ex]{\textwidth}{2pt}

\begin{questions}

%\question 
%
%\begin{parts}
%\part[2] 
%\begin{solution}
%\end{solution}
%
%
%\end{parts}
%\addpoints

\question[2] Determine los valores de $a$ y $b$ para sea
continua la función:
$$f(x)=\left\{\begin{array}{ll}
\dfrac{1}{e^x} & si \ x\leq 0 \\
a\cos(x) + b & si \  0 < x \leq \pi \\
\sen(x)-ax & si \ \pi<x
\end{array}\right.$$
\begin{solution}
Fuera de $x=0$ y $x=\pi$ la función es continua por serlo sus trozos en sus dominios. \\

En $x=0$: \\
$\lim_{x \to 0^-}f=\frac{1}{e^0}=1$\\
$\lim_{x \to 0^+}f=a\cdot \cos(0) +b =a+b \to a+b=1$\\

En $x=\pi$: \\
$\lim_{x \to \pi^-}f=a\cdot \cos(\pi) +b =b-a$\\
$\lim_{x \to \pi^+}f=\sen(\pi) - a\cdot \pi =- a\cdot \pi \to b-a=-a\cdot \pi$\\

$(2-\pi)\cdot a= 1 \to a=\frac{1}{2- \pi} \to b=\frac{1-\pi}{2-\pi}$

\end{solution}

\question Dada la función $f(x)= x^6-4x^4$ se pide:
\begin{parts}
    \part[1] Estudiar los intervalos de crecimiento y decrecimiento
    \begin{solution}
        $f=x^{6} - 4 x^{4} \to f'(x)=6 x^{5} - 16 x^{3} \to 30 x^{4} - 48 x^{2}$\\
        Puntos críticos: $f'(x)=0 \to  x=0, \ x= - \frac{2 \sqrt{6}}{3}, \  x=\frac{2 \sqrt{6}}{3} \to dec, crec, dec, crec $
    \end{solution}
    \part[1] Encontrar los máximos y mínimos relativos y determinar si son o no absolutos
    \begin{solution}
        Relativos: $x=0 \ max, \  x= - \frac{2 \sqrt{6}}{3} \ min, \  x=\frac{2 \sqrt{6}}{3} \ min$
        Absolutos: Max = $\infty$, Min = $x= - \frac{2 \sqrt{6}}{3} \ , \  x=\frac{2 \sqrt{6}}{3}$
    \end{solution}
    \part[1] Encontrar los puntos de inflexión
    \begin{solution}
        $f''(x)=0 \land f'''(x) \neq 0 \to  \  x=- \frac{2 \sqrt{10}}{5}, \  x=\frac{2 \sqrt{10}}{5}$
    \end{solution}
\end{parts}

\question Calcula las siguientes integrales:
\begin{parts}
\part[1]
$$\int x \ln{\left(x \right)}\, dx$$
\begin{solution}
$\left|\begin{matrix}
u = \ln x & du=\frac{1}{x} \, dx \\ dv= x \, dx & v=\frac{x^3}{3}
\end{matrix}\right| \to \frac{x^{2} \log{\left(x \right)}}{2} - \frac{x^{2}}{4}+K$
\end{solution}

\part[2]
$$\int x^{2} \ln^2{\left(x \right)}\, dx$$
\begin{solution}
$\left|\begin{matrix}
u = \ln^2 x & du=\frac{2}{x}\cdot \ln x \, dx \\ dv= x^2 \, dx & v=\frac{x^3}{3}
\end{matrix}\right| \to \frac{x^{3} \ln{\left(x \right)}^{2}}{3} - \frac{2 x^{3} \ln{\left(x \right)}}{9} + \frac{2 x^{3}}{27}+K$
\end{solution}

\part[2]
$$\int\limits_{2}^{3} \frac{1}{2 x^{2} - 4 x + 2}\, dx$$
\begin{solution}
$\frac{1}{4}  \land F(x)= -\frac{1}{2 x - 2}$
\end{solution}
\end{parts}


% \question[3] Calcule las dimensiones de tres campos cuadrados que no tienen ningún lado común y que
% satisfacen que el perímetro de uno de ellos es triple que el de otro y, además, se necesitan 1248 metros de
% valla para vallar completamente los tres campos, de manera que la suma de las áreas es la mínima posible.
% \begin{solution}
%     Lados= $x$, $3x$ e $y$ \\
%     $4x+4(3x)+4y=1248 \to 4(x+3x+y)=1248 \to 4x+y=312 \to y=312-4x$ \\
%     maximizar $f(x)=x^2+(3x)^2+(312-4x)^2=26 x^{2} - 2496 x + 97344= x^{2} - 96 x + 3744$ \\
%     $f'(x)=0 \to 2x-96=0 \to x=48 \land f''(48)=2>0 \to$ en $x=48$ hay un mínimo.
%     Solución: 48, 144, 120 m  
    
% \end{solution}

\question[2] Descomponer el número 12 en dos sumandos positivos de forma que el producto
del primero por el cuadrado del segundo sea máximo.
\begin{solution}
    $x$ e $y \to x+y =12 \to y=12-x$ \\
    Maximizar $f(x)=x\cdot y^2 \to f(x)=x\left(12-x\right)^2$ \\
    $f(x)=x^{3} - 24 x^{2} + 144 x \to f'(x)=3 x^{2} - 48 x + 144$ \\
    $f'(x)=0 \to x=4 y x=12 \to f''(x)=6x-48 \to f''(4)=-24<0 \land f''(12)=24>0 \to$ Solución. $x=4, y=8$
\end{solution}

\question[2] Calcule el valor de $a \in \mathbb{R}$, con $a\neq 0$, para que: $$\lim_{x \to 0}\dfrac{e^{a\cdot x^2}-\cos{x}}{\ln\left({x^2+1}\right)}=\frac{5}{6}$$
\begin{solution}
    L'H $\lim = \frac{2a+1}{2} \to a = \frac{1}{3}$
\end{solution}

% \question[2] Calcule el valor de $k$ para que: $$\lim_{x \to 0}\dfrac{e^x-e^{-x}+kx}{x-\sen{x}}=2$$
% \begin{solution}
%     L'H $\lim =\frac{2+k}{0}  \to k = -2$
% \end{solution}
% \addpoints

% % problema 23.2.2 p523
% \question[2] Determine el área de la región limitada inferiormente por la parábola $y=9x-x^2$ y superiormente por las rectas tangentes a la parábola en los puntos de corte de ella con el eje $OX$



% problema 3.3.2 p69
\question Sea la función $f(x)=x\cdot e^{x^2}$
\begin{parts}
    \part[1]Hallar la ecuación de la recta tangente a la gráfica de $f (x)$ en el punto de abscisa $x = 1$.
    \begin{solution}
        $y=e\left(3x-2\right)$
    \end{solution}
    \part[1] Calcula el área del recinto plano acotado limitado por la gráfica de $f(x)$ para $x \geq 0$, el eje $OX$ y la recta $x = 2$.
    \begin{solution}
        $\frac{e^4}{2}-\frac{1}{2} \ u^2$ 
    \end{solution}
\end{parts}
\addpoints



\end{questions}

\end{document}
\grid
